\chapter{Esercizi: Gestione delle eccezioni e dei file in Java}
\label{chap:esercizi-eccezioni-file-java}

\begin{exercise}
\label{ex:fi-esercizi-java-file-passaggio-valore}
Che cosa significa dire che in Java il passaggio dei parametri è
sempre per valore?
\end{exercise}

Il metodo riceve una copia del valore dell'argomento. Per un tipo
primitivo la copia contiene il dato stesso; per un tipo riferimento
contiene un riferimento allo stesso oggetto. Modificare un parametro
non modifica la variabile del chiamante, ma attraverso un riferimento
copiato si può modificare lo stato dell'oggetto condiviso
(\cref{def:fi-java-file-passaggio-valore}).

\begin{exercise}
\label{ex:fi-esercizi-java-file-metodo-non-statico}
Perché un metodo non statico come \texttt{deposit()} non può essere
invocato direttamente da \texttt{main()} senza indicare l'oggetto
su cui applicarlo?
\end{exercise}

\texttt{main()} è statico e non possiede un ricevente implicito
\texttt{this}. Il metodo \texttt{deposit()} opera invece sullo
stato di un oggetto specifico. Si deve quindi indicarne il
riferimento, come in \texttt{account.deposit(500)}
(\cref{sec:fi-java-file-metodi-ricevente}).

\begin{exercise}
\label{ex:fi-esercizi-java-file-this}
Che cosa significa dire che \texttt{this} è il parametro implicito
di un metodo non statico?
\end{exercise}

\texttt{this} individua l'oggetto ricevente sul quale è stata
invocata l'operazione. Consente di accedere ai campi e ai metodi
di istanza di quello stesso oggetto, anche quando il riferimento
non compare esplicitamente davanti al nome del membro
(\cref{sec:fi-java-file-metodi-ricevente}).

\begin{exercise}
\label{ex:fi-esercizi-java-file-costruttore-richiamo}
Perché non è possibile richiamare un costruttore su un oggetto
già esistente scrivendo \texttt{account.BankAccount()}?
\end{exercise}

Un costruttore è eseguito durante la creazione di un'istanza e
non è un metodo ordinario invocabile su un oggetto esistente.
L'espressione indicata tenta di individuare un metodo denominato
\texttt{BankAccount}, che non è definito
(\cref{sec:fi-java-file-ricostruire}).

\begin{exercise}
\label{ex:fi-esercizi-java-file-ricostruire}
Qual è il modo corretto per riportare una variabile oggetto a un
nuovo stato iniziale?
\end{exercise}

Si può creare un nuovo oggetto e assegnarne il riferimento alla
variabile, per esempio con \texttt{account = new BankAccount();}.
L'oggetto precedentemente indicato non viene azzerato: resta
raggiungibile se ne esistono altri riferimenti
(\cref{sec:fi-java-file-ricostruire}).

\begin{exercise}
\label{ex:fi-esercizi-java-file-package}
A che cosa servono i package in Java?
\end{exercise}

I package raggruppano tipi correlati in spazi dei nomi distinti,
facilitando l'organizzazione di programmi composti da più classi
e riducendo i conflitti fra nomi. Contribuiscono anche alle
regole di accessibilità
(\cref{def:fi-java-file-package}).

\begin{exercise}
\label{ex:fi-esercizi-java-file-package-default}
Che cos'è il package default e perché non è consigliato nei
progetti reali?
\end{exercise}

È il package senza nome a cui appartengono le classi prive di una
dichiarazione \texttt{package}. È adatto a esempi isolati ma
ostacola la modularità: le classi appartenenti a package con nome
non possono importarne i tipi
(\cref{sec:fi-java-file-package-struttura}).

\clearpage

\begin{exercise}
\label{ex:fi-esercizi-java-file-opzione-d}
A che cosa serve l'opzione \texttt{-d} del comando \texttt{javac}?
\end{exercise}

Specifica la directory di destinazione dei file \texttt{.class}
compilati. Il compilatore vi organizza i risultati nella struttura
corrispondente ai package; ciò consente di separare i sorgenti
dai prodotti della compilazione
(\cref{sec:fi-java-file-package-compilazione}).

\begin{exercise}
\label{ex:fi-esercizi-java-file-precondizioni}
Che cosa sono le precondizioni di un metodo e perché il
compilatore non può sempre verificarle?
\end{exercise}

Sono condizioni che gli argomenti e lo stato devono soddisfare
prima dell'operazione perché ne sia garantito il comportamento.
Il compilatore controlla le regole statiche di tipo, ma non
può stabilire in generale se, durante l'esecuzione, un importo
sia positivo o un indice appartenga all'intervallo ammesso
(\cref{def:fi-java-file-precondizione}).

\begin{exercise}
\label{ex:fi-esercizi-java-file-argomenti-inattesi}
Quali sono i possibili modi per reagire a un argomento inatteso,
e perché lanciare un'eccezione è spesso la soluzione migliore?
\end{exercise}

Il metodo potrebbe ignorare il valore, terminare il programma,
verificare un'asserzione oppure lanciare un'eccezione. Ignorare
l'errore può nasconderlo; terminare il programma impedisce al
chiamante di recuperare; un'asserzione può essere disabilitata.
L'eccezione segnala invece esplicitamente la condizione anomala
e lascia al chiamante la possibilità di gestirla
(\cref{sec:fi-java-file-precondizioni}).

\begin{exercise}
\label{ex:fi-esercizi-java-file-assert}
A che cosa serve un'asserzione in Java e perché non deve
sostituire la gestione delle eccezioni?
\end{exercise}

\texttt{assert} verifica condizioni che il programmatore ritiene
debbano essere vere, prevalentemente per individuare errori
interni durante sviluppo e collaudo. Le asserzioni sono normalmente
disabilitate e non possono quindi garantire la validazione degli
input né il recupero dagli errori
(\cref{sec:fi-java-file-precondizioni}).

\begin{exercise}
\label{ex:fi-esercizi-java-file-throw-propagazione}
Che cosa succede quando un metodo esegue \texttt{throw} e
l'eccezione non viene catturata nello stesso metodo?
\end{exercise}

La normale esecuzione del metodo viene interrotta; l'eccezione
si propaga al metodo chiamante e, se non viene intercettata,
prosegue lungo la catena delle chiamate. Un'eccezione non gestita
nel thread principale conduce normalmente alla sua terminazione
anomala con diagnostica
(\cref{def:fi-java-file-eccezione}).

\begin{exercise}
\label{ex:fi-esercizi-java-file-checked-unchecked}
Qual è la differenza tra eccezioni controllate e non
controllate in Java?
\end{exercise}

Per le eccezioni controllate il compilatore impone la gestione
mediante \texttt{catch} oppure la dichiarazione di propagazione
mediante \texttt{throws}. Questa prescrizione non vale per
\texttt{RuntimeException} e \texttt{Error}, né per le loro
sottoclassi
(\cref{def:fi-java-file-checked}).

\begin{exercise}
\label{ex:fi-esercizi-java-file-io-checked}
Perché \texttt{IOException} è un esempio tipico di eccezione
controllata?
\end{exercise}

Rappresenta fallimenti di ingresso/uscita, come problemi di
lettura, di scrittura o di accesso a risorse esterne, che possono
verificarsi anche quando il programma rispetta correttamente
la propria logica. Il compilatore richiede dunque che un metodo
capace di generarla la gestisca o ne dichiari la propagazione
(\cref{sec:fi-java-file-io-throws}).

\clearpage

\begin{exercise}
\label{ex:fi-esercizi-java-file-runtime-unchecked}
Perché \texttt{RuntimeException} e sue sottoclassi, come
\texttt{NumberFormatException}, sono eccezioni non controllate?
\end{exercise}

Java non impone al compilatore di verificare che siano
intercettate o dichiarate. Una conversione non valida può quindi
propagarsi senza una clausola \texttt{throws}; ciò non impedisce
al programmatore di intercettarla quando può gestire il dato
non valido
(\cref{def:fi-java-file-checked,sec:fi-java-file-throw-try-catch}).

\begin{exercise}
\label{ex:fi-esercizi-java-file-try}
A che cosa serve un blocco \texttt{try}?
\end{exercise}

Delimita le istruzioni eseguite in un contesto in cui è possibile
intercettare eccezioni mediante clausole \texttt{catch}, oppure
eseguire una chiusura o un'operazione di ripristino mediante
\texttt{finally}. Se viene lanciata un'eccezione, le istruzioni
successive dello stesso blocco non vengono eseguite normalmente
(\cref{def:fi-java-file-try-catch}).

\begin{exercise}
\label{ex:fi-esercizi-java-file-catch}
A che cosa serve una clausola \texttt{catch}?
\end{exercise}

Intercetta un'eccezione compatibile con il tipo dichiarato e
ne esegue il gestore. Permette, per esempio, di segnalare un
input non valido e riprovare, evitando che l'eccezione si
propaghi senza controllo
(\cref{sec:fi-java-file-throw-try-catch}).

\begin{exercise}
\label{ex:fi-esercizi-java-file-lanciare-catturare}
Che cosa significa il principio «lanciare presto, catturare tardi»?
\end{exercise}

Un metodo deve segnalare un problema quando lo rileva se non può
risolverlo correttamente. L'eccezione va intercettata soltanto
a un livello che disponga delle informazioni per recuperare
o comunicare il fallimento senza nasconderlo
(\cref{sec:fi-java-file-throw-try-catch}).

\begin{exercise}
\label{ex:fi-esercizi-java-file-parseint}
Perché un errore rilevato in \texttt{Integer.parseInt()} può
dover essere gestito in un punto diverso del programma?
\end{exercise}

Il metodo di conversione può stabilire che la stringa non
rappresenta un intero, ma non conosce il significato di quel dato
nel programma chiamante. Quest'ultimo può decidere se chiedere un
nuovo valore, segnalare un errore all'utente o interrompere
un'elaborazione
(\cref{sec:fi-java-file-throw-try-catch}).

\begin{exercise}
\label{ex:fi-esercizi-java-file-close}
Perché i flussi di ingresso e uscita devono essere chiusi dopo
l'uso, e che cosa può accadere se un'eccezione interviene
prima dell'invocazione di \texttt{close()}?
\end{exercise}

La chiusura libera risorse del sistema operativo e completa
le eventuali scritture ancora presenti nei buffer. Se un'eccezione
impedisce di eseguire \texttt{close()}, il programma può lasciare
risorse aperte o dati non scritti. \texttt{try-with-resources}
esegue automaticamente la chiusura anche quando un'operazione
fallisce
(\cref{sec:fi-java-file-try-resources}).

\begin{exercise}
\label{ex:fi-esercizi-java-file-sentinella-eof}
Qual è la differenza tra terminare l'input con un valore
sentinella e terminarlo comunicando al sistema operativo la
fine del flusso, per esempio con \texttt{Ctrl+D} o
\texttt{Ctrl+Z}?
\end{exercise}

La sentinella è un dato ordinario cui il programma attribuisce
il significato convenzionale di fine acquisizione. La fine del
flusso, invece, è una condizione segnalata dal sistema di I/O;
non richiede l'inserimento di un token speciale. Il metodo
\texttt{hasNextLine()} permette di rilevarla, restituendo
\texttt{false}
(\cref{sec:fi-java-file-eof-token}).
